% Generated draft: requires mathematical and source review.
\documentclass[11pt]{article}
\usepackage[margin=25mm]{geometry}
\usepackage{fontspec}
\setmainfont{Latin Modern Roman}
\usepackage{amsmath,amssymb,mathtools,unicode-math}
\setmathfont{Latin Modern Math}
\usepackage{enumitem,needspace,xurl,hyperref}
\setlength{\parindent}{0pt}
\setlength{\parskip}{5pt}
\setlength{\emergencystretch}{3em}
\newcommand{\N}{\mathbb{N}}
\newcommand{\Z}{\mathbb{Z}}
\newcommand{\Q}{\mathbb{Q}}
\newcommand{\R}{\mathbb{R}}
\newcommand{\C}{\mathbb{C}}
\newcommand{\F}{\mathbb{F}}
\newcommand{\A}{\mathbb{A}}
\newcommand{\PP}{\mathbb{P}}
\newcommand{\E}{\mathbb{E}}
\newcommand{\eps}{\varepsilon}
\newcommand{\dd}{\,\mathrm{d}}
\begin{document}
\section*{UnsolvedMath ID: 30006398}
\textbf{Source problem code:} catalogue entry
\subsection{Definitions and mathematical statement}
Let \(k\geq 3\) be a fixed integer, and let \(r\) denote the number of colors, interpreted as a positive integer. The supplied record calls \(P_r(k)\) a color-pattern parameter and \(f_{k,k+1}(r)\) the Erd\H{o}s--Rogers function. It does not supply definitions of these quantities. In particular, it does not specify the extremal conditions defining \(P_r(k)\). Their exact definitions must therefore be recovered from the cited source before this can be a self-contained mathematical statement.

With these symbols retaining their intended source meanings, the extracted conjecture is
\[
\text{for every fixed integer }k\geq 3,\qquad
P_r(k)=\Theta\!\left(r\,f_{k,k+1}(r)^2\right)
\quad\text{as }r\to\infty.
\]
Explicitly, the intended order-of-growth assertion is that for every integer \(k\geq 3\), there exist constants \(c_k,C_k>0\) and an integer \(r_0(k)\geq 1\) such that every integer \(r\geq r_0(k)\) satisfies
\[
c_k\,r\,f_{k,k+1}(r)^2
\leq P_r(k)
\leq C_k\,r\,f_{k,k+1}(r)^2.
\]
The constants and threshold may depend on \(k\); no uniform assertion for varying \(k\) is made. The limit through positive integers is the natural interpretation of the color-count parameter, but the excerpt does not explicitly state the limit or its domain.

The supplied dated curation labels the problem partially solved and says that a construction comes within a logarithmic factor of the conjectured order, while the full asymptotic order remains unproved. This is a supplied literature assessment, not an independently verified conclusion. The excerpt refers to an unspecified Theorem 2 and to improvements of Mubayi and Verstraete, but supplies neither the theorem nor precise bounds; no more specific partial result is asserted here.
\subsection{Short English statement}
For each fixed integer \(k\geq 3\), the conjecture asks whether the color-pattern parameter \(P_r(k)\) has order \(r\,f_{k,k+1}(r)^2\) as the number of colors \(r\) grows. The comparison constants may depend on \(k\). The supplied record omits the definitions of both parameters, so their precise source conventions require review. Its partially solved label and logarithmic-factor assessment are supplied dated curation.
\subsection{Sources}
\begin{itemize}
\item Local UnsolvedMath registry: \texttt{data/unsolvedmath/problems.json}, permanent ID \(30006398\), source problem code \texttt{catalogue entry}, title ``Color Patterns and the Erd\H{o}s--Rogers Function.'' This draft uses the supplied complete registry record.
\item Supplied report citation: \emph{Combinatorics, Probability and Computing (2026)}, Oberwolfach Reports, DOI 10.4171/owr/2025/42, pp.~2243--2322. \url{https://doi.org/10.4171/owr/2025/42}. The record separately assigns the proposed year 2025; the year discrepancy is unresolved.
\item Supplied literature source: ``Oberwolfach Report 42 2025.'' \url{https://oa.tib.eu/renate/bitstreams/72e281a6-ba33-4ac7-baf7-6fadd1126345/download}.
\item Supplied literature source: \emph{Improved Bounds for Minimal Multicolor Ramsey Graphs}. \url{https://arxiv.org/abs/2510.09068}. The supplied record gives no authors for this reference.
\item Literature assessment dated exactly \texttt{2026-08-22}: the record labels the problem partially solved and describes bounds within a logarithmic factor of the conjectured order. This date and assessment are supplied dated curation, not independently verified source or literature checks.
\item Attribution limitations: the excerpt refers to ``Conjecture 5.2'' without its bibliographic reference, to Mubayi and Verstraete [9], and to constructions of Fox et al.\ [6], Dudek, Retter, and R\"odl [4], H\`an et al.\ [8], and Bamberg et al.\ [1]. The numbered bibliography is absent. No individual proposer is identified in the supplied record.
\end{itemize}
\end{document}
